% Standalone resolution manuscript generated from solution.md.
% Compile with XeLaTeX; author and review metadata are maintained in Markdown.
\documentclass[10pt,a4paper]{article}
\usepackage[margin=24mm,headheight=15pt]{geometry}
\usepackage{fontspec}
\setmainfont{DejaVuSerif.ttf}[BoldFont=DejaVuSerif-Bold.ttf,ItalicFont=DejaVuSerif-Italic.ttf,BoldItalicFont=DejaVuSerif-BoldItalic.ttf]
\setsansfont{DejaVuSans.ttf}[BoldFont=DejaVuSans-Bold.ttf,ItalicFont=DejaVuSans-Oblique.ttf]
\setmonofont{DejaVuSansMono.ttf}
\usepackage{amsmath,amssymb,unicode-math}
\setmathfont{latinmodern-math.otf}
\usepackage{xcolor,microtype,parskip,needspace,fancyhdr}
\usepackage{bookmark}
\usepackage{xurl}
\hypersetup{colorlinks=true,urlcolor=blue,linkcolor=blue,pdftitle={MI-24:
the Schatten norm complement follows by convexity},pdfauthor={George
Stepaniants},pdflang={en-GB}}
\urlstyle{same}
\setlength{\emergencystretch}{3em}
\providecommand{\tightlist}{\setlength{\itemsep}{0pt}\setlength{\parskip}{0pt}}
\providecommand{\pandocbounded}[1]{#1}
\setcounter{secnumdepth}{0}
\allowdisplaybreaks[1]
\widowpenalty=10000
\clubpenalty=10000
\pagestyle{fancy}
\fancyhf{}
\fancyhead[L]{\small\sffamily Verified resolution / MI-24}
\fancyhead[R]{\small\sffamily 11 September 2026}
\fancyfoot[L]{\parbox[b]{0.9\textwidth}{\fontsize{8}{10}\selectfont Independent
Codex-agent review; not external human peer review or formal
certification.}}
\fancyfoot[R]{\footnotesize\thepage}
\begin{document}
{\raggedright\Large\bfseries MI-24: the Schatten norm complement follows
by convexity\par}
\vspace{0.4em}
{\normalsize\bfseries George Stepaniants\par}
{\small Department of Computing and Mathematical Sciences, California
Institute of Technology, Pasadena, California, USA\par}
\vspace{0.6em}
\pagestyle{plain}

The full Schatten norm inequality in MI-24 follows from the published
Heron norm inequality of Dinh, Dumitru, and Franco, a positive matrix
comparison already used by Ghabries, Abbas, Mourad, and Assi, and the
triangle inequality. The last step is a general convexity observation:
if one endpoint of a segment has norm at most the midpoint's norm, then
the other endpoint has norm at least the midpoint's norm.

\textbf{Verification status: independent-agent PASS.} Developed with
substantial ChatGPT/Codex assistance. A separate agent checked the
complete argument and its primary-source application.
\href{https://github.com/ajt60gaibb/OpenProblemsInNLA/blob/main/references/stepaniants-mi24-2026-09-11/verification/reviews/MI-24-review.md}{Detailed
independent review}. This is independent automated review, not external
human peer review or formal verification.

\subsection{The full statement}\label{the-full-statement}

Let \(A,B\in\mathbb C^{n\times n}\) be positive definite, where
\(n\ge1\), and define

\[
\begin{aligned}
G&=A^{1/2}(A^{-1/2}BA^{-1/2})^{1/2}A^{1/2},\\
L&=A^{1/2}(B^{1/2}A^{-1}B^{1/2})^{1/2}A^{1/2}.
\end{aligned}
\]

All square roots are the positive definite square roots. For
\(1\le p<\infty\), write \(\|M\|_p=(\operatorname{tr}|M|^p)^{1/p}\),
where \(|M|=(M^*M)^{1/2}\); \(\|M\|_\infty\) is the largest singular
value.

\subsubsection{Theorem 1}\label{theorem-1}

For every such pair and every \(1\le p\le\infty\),

\[
\boxed{\|A+B+G+L\|_p\le\|A+B+2L\|_p.}
\]

This resolves the complete target of
\href{https://github.com/ajt60gaibb/OpenProblemsInNLA/blob/main/matrix-inequalities-and-norms/MI-24/README.md}{MI-24},
equivalently
\href{https://arxiv.org/html/2105.13356v1\#S4}{Ghabries--Abbas--Mourad--Assi,
Conjecture 4.1}, including the previously unproved exponents. Neither
commutativity nor symmetry of \(L(A,B)\) in its arguments is assumed.

\subsection{The two comparisons}\label{the-two-comparisons}

Put \(H=A+B\) and \(K=A^{1/2}B^{1/2}+B^{1/2}A^{1/2}\). The published
Heron norm inequality of
\href{https://trunghoamath.wordpress.com/wp-content/uploads/2017/06/p_t_q_t-4-2.pdf}{Dinh--Dumitru--Franco,
Theorem 4}, p.~2 of the author manuscript, gives

\[
\|H+2G\|_p\le\|H+K\|_p \qquad(1\le p<\infty). \tag{1}
\]

The same inequality at \(p=\infty\) follows by letting \(p\to\infty\) in
finite dimension. This is the only non-elementary external theorem
needed below; it is also recalled in equation (22) and its following
paragraph in
\href{https://arxiv.org/html/2105.13356v1\#S4}{Ghabries--Abbas--Mourad--Assi}.

The second comparison is

\[
K\preceq G+L. \tag{2}
\]

For completeness, it has the following direct proof. Set
\(R=A^{-1/2}B^{1/2}\) and take its polar decomposition \(R=U|R|\), where
\(U\) is unitary. Since \(|R^*|=U|R|U^*\),

\[
|R^*|+|R|-R-R^*=(U-I)|R|(U^*-I)\succeq0.
\]

Congruence by \(A^{1/2}\) gives (2), because

\[
\begin{aligned}
A^{1/2}|R^*|A^{1/2}&=G,&
A^{1/2}|R|A^{1/2}&=L,\\
A^{1/2}(R+R^*)A^{1/2}&=K.
\end{aligned}
\]

Moreover, \(H+K=(A^{1/2}+B^{1/2})^2\succeq0\). Hence (2) implies

\[
0\preceq H+K\preceq H+G+L.
\]

The min--max principle orders every eigenvalue of these positive
semidefinite matrices, so

\[
\|H+K\|_p\le\|H+G+L\|_p \qquad(1\le p\le\infty). \tag{3}
\]

This also reproduces the particular comparison in
Ghabries--Abbas--Mourad--Assi, Corollary 4.1, without requiring their
block-matrix proof as an additional premise.

\subsection{The convexity step}\label{the-convexity-step}

\textbf{Proof of Theorem 1.} Fix any \(p\in[1,\infty]\) and set

\[
X=H+2G,\qquad Y=H+G+L,\qquad Z=H+2L.
\]

Equations (1) and (3) give \(\|X\|_p\le\|Y\|_p\). Also \(X+Z=2Y\). The
triangle inequality therefore gives

\[
2\|Y\|_p=\|X+Z\|_p
\le\|X\|_p+\|Z\|_p
\le\|Y\|_p+\|Z\|_p.
\]

Subtracting \(\|Y\|_p\) proves \(\|Y\|_p\le\|Z\|_p\), as claimed.
\(\square\)

The proof covers all parameters and dimensions in the canonical
statement. Its contribution is the convexity deduction connecting the
two existing comparisons; no new proof of the published Heron inequality
or historical priority certification is claimed.

\subsection{References}\label{references}

\begin{enumerate}
\def\labelenumi{\arabic{enumi}.}
\item
  \textbf{Dinh--Dumitru--Franco (2017).} T. H. Dinh, R. Dumitru, and J.
  A. Franco. \emph{On a conjecture of Bhatia, Lim and Yamazaki}. Linear
  Algebra and its Applications \textbf{532} (2017), 140--145. Theorem 4,
  equation (4), p.~2 of the
  \href{https://trunghoamath.wordpress.com/wp-content/uploads/2017/06/p_t_q_t-4-2.pdf}{author
  manuscript}. \href{https://doi.org/10.1016/j.laa.2017.06.040}{DOI:
  10.1016/j.laa.2017.06.040}.
\item
  \textbf{Ghabries--Abbas--Mourad--Assi (2021/2022).} M. M. Ghabries, H.
  Abbas, B. Mourad, and A. Assi. \emph{New log-majorization results
  concerning eigenvalues and singular values and a complement of a norm
  inequality}. Section 4: equation (22), Corollary 4.1, Theorem 4.2 and
  Conjecture 4.1, pp.~11--13 of
  \href{https://arxiv.org/abs/2105.13356v1}{arXiv:2105.13356v1}.
  \href{https://doi.org/10.1080/03081087.2022.2059050}{Journal DOI:
  10.1080/03081087.2022.2059050}.
\item
  \textbf{MI-24.} ajt60gaibb/OpenProblemsInNLA. \emph{Schatten norm
  complement involving Lin's positive definite quantity}.
  \href{https://github.com/ajt60gaibb/OpenProblemsInNLA/blob/main/matrix-inequalities-and-norms/MI-24/README.md}{Canonical
  problem statement}. The original mathematical target and permanent
  identifier are retained.
\end{enumerate}
\end{document}
